\chapter{Saturation flow and lane capacity}
\label{chap:saturation-flow}

\section{Definition}
Saturation flow is the maximum queue-discharge flow rate at a stop line under continuous
effective service, expressed in PCU/h of effective green.  It is not the same as lane
capacity over a complete signal cycle.  Nominal capacity is obtained by applying the
effective-green ratio.

\section{Direct input}
A measured or otherwise justified saturation flow may be entered directly:
\begin{verbatim}
saturation_flow:
  method: direct
  value_pcu_per_hour: 1900
\end{verbatim}
Field measurement is preferable where suitable observations are available because
saturation flow strongly affects capacity, DoS, queue and delay.

\section{RR67-style estimation}
PySig also implements an RR67-style unopposed saturation-flow estimate.  Supported
inputs include lane width, gradient, nearside-lane status, turning proportion,
turning radius and a local calibration factor.  The software output should expose both
the final saturation-flow value and the inputs used to obtain it.

Example:
\begin{verbatim}
saturation_flow:
  method: rr67_unopposed
  lane_width_m: 3.2
  gradient_percent: 1.5
  is_nearside_lane: false
  turning_radius_m: 18.0
  local_factor: 1.0
\end{verbatim}

\section{Dynamic turning proportion}
If neither \code{turning\_proportion} nor \code{turning\_proportion\_override} is supplied,
PySig derives the current lane turning proportion from assignment.  This avoids treating
a behavioural lane-flow property as permanently fixed before assignment.

A historical \code{turning\_proportion} remains a fixed-input override for compatibility.
The preferred explicit form for a calibrated fixed assumption is
\code{turning\_proportion\_override}.

\section{Turning radius}
Turning radius is supplied in metres.  Where one turn path dominates the lane,
\code{turning\_radius\_m} may be stored in the lane saturation-flow block.  Where several
turning movements provide movement-specific radii, PySig can derive a turning-flow-
weighted radius.  This is an assignment-responsive approximation rather than a claim
that the weighted radius has independent geometric meaning.

\section{Base and accessible capacity}
RR67 returns a base saturation-flow estimate $s$.  Effective green gives nominal
capacity
\begin{equation}
c_n=s\frac{g}{C}.
\end{equation}
Short lanes, shared-lane blocking and priority effects can then reduce or redistribute
the service accessible to individual lanes.  For auditability, these interactions do
not overwrite the base RR67 value.

\section{Calibration}
A calculated saturation flow is a prediction.  Site-specific measured discharge should
be used to calibrate or override prediction where adequate data exist.  Queue discharge,
start loss/end gain and saturation flow are observable field quantities and should be
checked together for important models.
